
\section{\textbf{Introdução}}
  Na astrofísica contemporânea, a determinação precisa da massa estelar se apresenta como um desafio fundamental. A massa de uma estrela exerce um papel central na evolução química das galáxias, na formação de planetas, na análise de sistemas estelares múltiplos e na interpretação dos períodos de rotação em aglomerados jovens. Dada a complexidade desses fenômenos, a necessidade de maior precisão nas estimativas de massa estelar se torna evidente.
  Em sistemas estelares múltiplos, a massa pode ser determinada diretamente. Para estrelas não-binárias na sequência principal, a estimativa de massa envolve frequentemente a relação massa-luminosidade, cujo expoente varia conforme o estágio evolutivo da estrela. Métodos de ajuste de isócronas se destacam para a determinação de massas de estrelas jovens de baixa massa, especialmente na fase pré-sequência principal. Esses métodos podem ser complementares, já que cada um possui vantagens e limitações específicas.
  Neste estudo, investigamos a determinação das massas individuais em aglomerados de estrelas jovens por meio de uma abordagem que integra dados de diferentes fontes, focado no Aglomerado das Plêiades. Utilizamos técnicas de aprendizado de máquina para estimar as massas estelares com base na relação massa-luminosidade, empregando dados de estrelas sintéticas provenientes de modelos teóricos de isócronas. Adicionalmente, aplicamos o método de ajuste de isócronas, voltado principalmente para aglomerados mais jovens.
  Os resultados obtidos foram comparados com os métodos tradicionais, incluindo o ajuste de isócronas e as relações empíricas e semi-empíricas da relação massa-luminosidade.